% GENERATED FILE. Edit canonical lessons, then run npm run generate:books.
% canonical-source-hash: fnv1a64:2bee5596b182162b
% canonical-lessons: AR-C21-kalb-qitt

\chapter{Dog and Cat}
\label{ch:ar-dog-and-cat}
\hlchaptermodality{21}

\begin{chapteropening}
\textbf{By the end of this chapter:} \emph{I can name a dog and a cat in Arabic.}

Say kalb beside Hebrew kelev, then qiṭṭ beside Latin cattus, flagging that the cat link is still debated.
\end{chapteropening}

\section[\texorpdfstring{\ar{كلب}, \ar{قط} (kalb, qiṭṭ)}{كلب, قط (kalb, qiṭṭ)}]{\ar{كلب}, \ar{قط} — a solid root, and the word that closes the loop}

\label{lesson:AR-C21-kalb-qitt}



\begin{quote}
Spanish has a genuine dog-word mystery, and German another by way of English. Arabic's dog-word has no such mystery — and its cat-word is tangled up, by some route, with nearly every European cat-word there is.
\end{quote}

\begin{culture}
\textbf{\ar{كلب}} (\emph{kalb}, \textquotedblleft{}\textbf{dog}\textquotedblright{}) traces to \textbf{Proto-Semitic} \emph{\textbf{*kalb-}} — the exact same root behind \textbf{Hebrew}'s own word for dog, \textbf{\he{כֶּלֶב}} (\emph{kelev}) (the name \textbf{Caleb} is traditionally connected to this same root, though some scholars instead derive it from \emph{kol-lev}, \textquotedblleft{}whole-hearted\textquotedblright{}). Unlike Spanish's \emph{perro} or the English word that replaced \textquotedblleft{}hound,\textquotedblright{} \emph{kalb} is a solid, well-attested, unmysterious Semitic root, shared across the whole Semitic family (Aramaic, Akkadian, and more).
\end{culture}

\begin{culture}
\textbf{\ar{قط}} (\emph{qiṭṭ}, \textquotedblleft{}\textbf{cat, tomcat}\textquotedblright{}) is widely \textbf{compared} with Latin's \textbf{cattus} (alongside Nubian \emph{kadis}, mentioned back in the Spanish and Latin lessons). Be honest about the uncertainty here, though: linguists don't agree on the relationship. \emph{Cattus} turns up late in Latin, and some think \emph{kadis} came \textbf{from} Arabic rather than the reverse. These look like one travelling word, spreading with the domestic cat, but where it started, and which language lent it to which, is not settled. Arabic's word for \textquotedblleft{}cat\textquotedblright{} sits in the same tangled word-family as Spanish's \emph{gato}, French's \emph{chat}, Italian's \emph{gatto}, and German's \emph{Katze}.
\end{culture}

\subsection*{Your turn}
Say these aloud:

\begin{itemize}
\raggedright
  \item \textquotedblleft{}kalb\textquotedblright{} — dog, a solid Semitic root, cousin of Hebrew's kelev
  \item \textquotedblleft{}qiṭṭ\textquotedblright{} — cat, linked by some route with Latin's cattus and every European cat-word there is
\end{itemize}

\begin{tcolorbox}[breakable,colback=teal!4,colframe=teal!35!black,title={Before you move on}]
What root does \textbf{\ar{كلب}} trace to, and what's its Hebrew cousin? (\textbf{Proto-Semitic \emph{*kalb-}}; cousin of Hebrew \textbf{\he{כֶּלֶב}}, \emph{kelev}.) Is \textbf{\ar{قط}}'s relationship to Latin's \emph{cattus} fully settled? (\textbf{No} — widely compared, probably one travelling word, but linguists still debate where it started and which way it went.) What does this most likely mean for Arabic's word for \textquotedblleft{}cat\textquotedblright{} compared to Spanish's \emph{gato}, French's \emph{chat}, and German's \emph{Katze}? (\textbf{They most likely belong to the same tangled ancient word family} — spreading with the domestic cat.)
\end{tcolorbox}
